\documentclass[../DoAn.tex]{subfiles}
\begin{document}

\section{Tổng quan chương}
Chương 3 đã trình bày các nền tảng lý thuyết và công nghệ được lựa chọn để xây dựng hệ thống Smart Grading. Trên cơ sở đó, Chương 4 trình bày chi tiết quá trình phân tích thiết kế, kết quả triển khai và đánh giá hệ thống. Nội dung chương gồm: thiết kế kiến trúc tổng thể và các gói thành phần, thiết kế chi tiết giao diện, lớp và cơ sở dữ liệu, kết quả xây dựng ứng dụng với minh họa các chức năng chính, kết quả kiểm thử và triển khai thử nghiệm trong môi trường thực tế.

\section{Thiết kế kiến trúc}
\label{section:4.1}

\subsection{Lựa chọn kiến trúc phần mềm}
\label{subsection:4.1.1}

Hệ thống Smart Grading lựa chọn kiến trúc phân lớp kết hợp với mô hình client--server, phù hợp với đặc thù của ứng dụng có nhiều thành phần tương tác qua mạng. Toàn bộ hệ thống được chia thành ba tầng chính:
\begin{itemize}
    \item tầng trình bày gồm ứng dụng di động Flutter và giao diện web React;
    \item tầng xử lý nghiệp vụ là backend Node.js/Express;
    \item tầng dữ liệu gồm MongoDB và dịch vụ lưu trữ đám mây Cloudinary.
\end{itemize}

Trong tầng backend, hệ thống áp dụng biến thể của mô hình MVC gồm bốn lớp: lớp định tuyến (\texttt{routes}) tiếp nhận và phân loại yêu cầu; lớp điều khiển (\texttt{controllers}) điều phối luồng xử lý và định dạng phản hồi; lớp dịch vụ (\texttt{services}) chứa toàn bộ logic nghiệp vụ; và lớp mô hình (\texttt{models}) quản lý tương tác với cơ sở dữ liệu. Các thành phần này được tổ chức thành các gói (module) độc lập theo từng miền nghiệp vụ như \texttt{auth}, \texttt{user}, \texttt{exam}, \texttt{submission}, \texttt{appeal}, \texttt{question} và \texttt{notification}.

Trong tầng giao diện web, kiến trúc áp dụng mô hình component với React và phân tách rõ ràng giữa ba lớp: lớp trình bày (pages và components), lớp quản lý trạng thái (Zustand stores) và lớp dịch vụ (services) để giao tiếp với backend thông qua REST API. Trong tầng di động, kiến trúc BLoC kết hợp với Repository pattern được áp dụng, trong đó BLoC xử lý logic nghiệp vụ, Repository đóng vai trò trung gian giữa BLoC và các nguồn dữ liệu (API server, lưu trữ cục bộ).

\subsection{Thiết kế tổng quan}
\label{subsection:4.1.2}

Biểu đồ gói UML của toàn bộ hệ thống được mô tả như trong Hình \ref{fig:package_diagram}, thể hiện sự phụ thuộc một chiều giữa các gói theo từng tầng.

\begin{figure}[H]
    \centering
    \includegraphics[width=0.9\textwidth]{Hinhve/Picture3.png}
    \caption{Biểu đồ gói tổng quan của hệ thống Smart Grading}
    \label{fig:package_diagram}
\end{figure}

\noindent\textbf{Nhóm gói tầng trình bày (Presentation Layer).} Gồm hai gói con: \texttt{web} chứa các trang và component React, \texttt{mobile} chứa các màn hình và widget Flutter. Các gói này chỉ phụ thuộc vào tầng dịch vụ chung và không phụ thuộc trực tiếp vào tầng dữ liệu.

\noindent\textbf{Nhóm gói tầng dịch vụ chung (Service Layer).} Gồm ba gói: \texttt{api-client} (web) và \texttt{network} (mobile) cung cấp các hàm tiện ích để gọi REST API; \texttt{state-management} quản lý trạng thái phía client; \texttt{routing} định tuyến và phân quyền truy cập theo vai trò.

\noindent\textbf{Nhóm gói tầng nghiệp vụ (Business Layer).} Gồm các gói backend: \texttt{routes} định nghĩa điểm cuối API; \texttt{controllers} xử lý request; \texttt{services} chứa logic nghiệp vụ được chia theo miền; \texttt{validations} xác thực đầu vào; \texttt{middlewares} xử lý xác thực, phân quyền và ghi log.

\noindent\textbf{Nhóm gói tầng dữ liệu (Data Layer).} Gồm các gói \texttt{models} (định nghĩa schema MongoDB thông qua Mongoose), \texttt{amc} (sinh và xử lý đề thi bằng AMC), \texttt{ai-services} (tích hợp các nhà cung cấp AI), \texttt{storage} (tương tác với Cloudinary), \texttt{utils} (các hàm tiện ích chung).

\subsection{Thiết kế chi tiết gói}
\label{subsection:4.1.3}

Biểu đồ lớp của nhóm gói nghiệp vụ chính (gồm các gói liên quan đến đề thi, bài làm và phúc khảo) được mô tả như trong Hình \ref{fig:class_diagram}, trong đó chỉ hiển thị tên lớp và quan hệ giữa các lớp.

\begin{figure}[H]
    \centering
    \includegraphics[width=0.9\textwidth]{Hinhve/Picture2.png}
    \caption{Ví dụ thiết kế gói biểu đồ lớp}
    \label{fig:class_diagram}
\end{figure}

Trong nhóm gói quản lý đề thi, các lớp chính gồm \texttt{ExamController}, \texttt{ExamService}, \texttt{ExamModel}, \texttt{ExamVersionModel}, \texttt{QuestionModel}, \texttt{OMRTemplateModel} và \texttt{ExamPaperService}. Lớp \texttt{ExamController} kế thừa \texttt{BaseController} và có quan hệ phụ thuộc với \texttt{ExamService}. Lớp \texttt{ExamService} có quan hệ kết tập với \texttt{ExamModel} và \texttt{ExamVersionModel}, đồng thời có quan hệ kết hợp với \texttt{ExamPaperService} để sinh đề thi bằng AMC. Lớp \texttt{ExamModel} thực thi giao diện \texttt{IDocument} của Mongoose.

Trong nhóm gói quản lý bài làm, các lớp chính gồm \texttt{SubmissionController}, \texttt{SubmissionService}, \texttt{SubmissionModel}, \texttt{AppealModel} và \texttt{AppealService}. Lớp \texttt{SubmissionModel} có quan hệ tham chiếu (reference) tới \texttt{ExamModel}, \texttt{ExamVersionModel} và \texttt{UserModel}. Khi một bài làm có phúc khảo, lớp \texttt{AppealService} sẽ tương tác với cả \texttt{SubmissionModel} và \texttt{ExamModel} để tính lại điểm khi phúc khảo được phê duyệt.

\section{Thiết kế chi tiết}
\label{section:4.2}

\subsection{Thiết kế giao diện}
\label{subsection:4.2.1}

Hệ thống Smart Grading hỗ trợ hai nhóm nền tảng giao diện với các chuẩn thiết kế thống nhất. Đối với ứng dụng di động Flutter, hệ thống hướng tới các thiết bị Android từ phiên bản 7.0 trở lên và iOS từ phiên bản 13 trở lên, kích thước màn hình tối thiểu 5.0 inch, độ phân giải tối thiểu HD (1280$\times$720). Đối với giao diện web, hệ thống hỗ trợ các trình duyệt phổ biến (Chrome, Edge, Firefox, Safari) ở độ phân giải tối thiểu 1366$\times$768 và tối ưu cho độ phân giải 1920$\times$1080.

Chuẩn phối màu được sử dụng xuyên suốt hệ thống gồm: màu chủ đạo xanh đậm \texttt{\#0b2240} cho tiêu đề và thanh điều hướng; màu phụ xanh dương \texttt{\#2563eb} cho các nút hành động chính; màu thành công xanh lá \texttt{\#059669} cho trạng thái tích cực; màu cảnh báo đỏ \texttt{\#dc2626} cho trạng thái lỗi; màu nền xám nhạt \texttt{\#f3f4f6} cho vùng nền. Phông chữ chính là phông không chân (sans-serif) để đảm bảo dễ đọc trên màn hình, với kích thước tiêu chuẩn 14pt cho nội dung, 16pt cho tiêu đề phụ và 20pt cho tiêu đề chính. Thông báo phản hồi được đặt ở góc trên bên phải màn hình với biểu tượng phù hợp, tự động biến mất sau 3 giây.

Hình \ref{fig:mockup_dashboard} và Hình \ref{fig:mockup_scan} minh họa một số màn hình chính của hệ thống.

\begin{figure}[H]
    \centering
    \includegraphics[width=0.85\textwidth]{Hinhve/Picture1.png}
    \caption{Minh họa mockup giao diện dashboard trên web}
    \label{fig:mockup_dashboard}
\end{figure}

\begin{figure}[H]
    \centering
    \includegraphics[width=0.85\textwidth]{Hinhve/Picture2.png}
    \caption{Minh họa mockup giao diện quét phiếu trên di động}
    \label{fig:mockup_scan}
\end{figure}

\subsection{Thiết kế lớp}
\label{subsection:4.2.2}

Trong phần này, đồ án trình bày chi tiết hai lớp chủ đạo của hệ thống là \texttt{ExamService} và \texttt{SubmissionService}, cùng với biểu đồ trình tự minh họa luồng truyền thông điệp cho hai use case quan trọng.

\noindent\textbf{Lớp ExamService.} Lớp này chịu trách nhiệm quản lý toàn bộ vòng đời của một đề thi từ khi tạo đến khi phát hành. Các thuộc tính chính gồm: \texttt{examModel} (tham chiếu đến model Exam), \texttt{versionModel} (tham chiếu đến model ExamVersion), \texttt{questionService} (tham chiếu đến service câu hỏi), \texttt{omrTemplateService} (tham chiếu đến service mẫu OMR), \texttt{examPaperService} (tham chiếu đến service sinh đề bằng AMC), \texttt{notificationService} (tham chiếu đến service thông báo), \texttt{aiService} (tham chiếu đến dịch vụ AI hỗ trợ sinh câu hỏi). Các phương thức chính gồm: \texttt{createExam(data)} tạo đề thi mới với danh sách câu hỏi và cấu hình trộn; \texttt{generateVersions(examId, numberOfVersions)} sinh các phiên bản đề bằng cách hoán vị câu hỏi và đáp án; \texttt{generatePapers(examId)} gọi service sinh PDF bằng AMC; \texttt{publishExam(examId)} chuyển trạng thái đề thi thành "đã phát hành" và gửi thông báo cho học sinh; \texttt{getExamById(examId)} truy vấn thông tin chi tiết đề thi.

\noindent\textbf{Lớp SubmissionService.} Lớp này chịu trách nhiệm xử lý bài làm của học sinh từ khi chấm đến khi công bố điểm và hỗ trợ phúc khảo. Các thuộc tính chính gồm: \texttt{submissionModel} (tham chiếu đến model Submission), \texttt{appealModel} (tham chiếu đến model Appeal), \texttt{examService} (tham chiếu đến service đề thi), \texttt{aiReportService} (tham chiếu đến service báo cáo AI). Các phương thức chính gồm: \texttt{createSubmissionFromScan(scanData)} tạo bài làm từ dữ liệu quét OMR; \texttt{manualOverride(submissionId, overrides)} cho phép giáo viên chỉnh sửa thủ công; \texttt{publishScores(examId)} công bố điểm cho học sinh; \texttt{recalculateScore(submissionId, questionId)} tính lại điểm khi phúc khảo được phê duyệt; \texttt{getSubmissionDetail(submissionId)} truy vấn chi tiết bài làm.

\noindent\textbf{Biểu đồ trình tự use case Chấm thi bằng ứng dụng.} Luồng truyền thông điệp cho use case Chấm thi bằng ứng dụng được mô tả tổng quan như sau. Đầu tiên, giáo viên mở chức năng quét phiếu trên ứng dụng di động; \texttt{MobileBloc} gửi yêu cầu lấy mẫu OMR đến \texttt{OMRTemplateService} trên server; server trả về mẫu OMR gồm tọa độ các ô; giáo viên chụp ảnh phiếu; module \texttt{OMREngine} trên di động thực hiện tiền xử lý ảnh và nhận dạng; kết quả nhận dạng được gửi đến \texttt{SubmissionController} thông qua \texttt{SubmissionService} để lưu vào cơ sở dữ liệu; hệ thống tính điểm và trả kết quả về ứng dụng di động; giáo viên kiểm tra và có thể chỉnh sửa; cuối cùng kết quả được lưu chính thức.

\noindent\textbf{Biểu đồ trình tự use case Phản hồi phúc khảo.} Luồng truyền thông điệp cho use case Phản hồi phúc khảo diễn ra như sau. Học sinh gửi đơn phúc khảo thông qua \texttt{AppealController} trên web; \texttt{AppealService} tạo đơn với trạng thái "đang chờ" và gửi thông báo cho giáo viên; giáo viên truy cập danh sách đơn, chọn đơn cần xử lý; \texttt{SubmissionService} trả về chi tiết bài làm; giáo viên quyết định phê duyệt hoặc từ chối; \texttt{AppealService} cập nhật trạng thái và gọi \texttt{SubmissionService.recalculateScore} nếu phê duyệt; hệ thống gửi thông báo kết quả cho học sinh.

\subsection{Thiết kế cơ sở dữ liệu}
\label{subsection:4.2.3}

Hệ thống sử dụng cơ sở dữ liệu MongoDB phiên bản 5.0 trở lên với 17 bảng (collection) chính, được thiết kế để đảm bảo tính nhất quán và hiệu quả truy vấn. Một số bảng cốt lõi gồm \texttt{users}, \texttt{schools}, \texttt{classes}, \texttt{subjects}, \texttt{questions}, \texttt{exams}, \texttt{examVersions}, \texttt{omrTemplates}, \texttt{submissions}, \texttt{appeals}, \texttt{notifications}, \texttt{aiReports}, \texttt{aiChats}, \texttt{studentProgress}, \texttt{examReports} và \texttt{uploadAuditLogs}.

\noindent\textbf{Bảng users.} Lưu trữ thông tin tài khoản người dùng với các trường chính: \texttt{\_id} (ObjectId), \texttt{name} (String), \texttt{email} (String, unique), \texttt{password} (String, hash), \texttt{role} (Enum: admin, school-admin, teacher, student, parent), \texttt{schoolId} (ObjectId tham chiếu schools), \texttt{classIds} (Array of ObjectId tham chiếu classes), \texttt{studentCode} (String, mã số sinh viên), \texttt{isEmailVerified} (Boolean), \texttt{timestamps} (createdAt, updatedAt).

\noindent\textbf{Bảng exams.} Lưu trữ thông tin đề thi với các trường chính: \texttt{\_id}, \texttt{title}, \texttt{description}, \texttt{classIds} (danh sách lớp áp dụng), \texttt{primaryClassId} (lớp chính), \texttt{omrTemplateId} (tham chiếu mẫu OMR), \texttt{questionIds} (danh sách câu hỏi), \texttt{examDate}, \texttt{duration} (phút), \texttt{totalScore}, \texttt{numberOfVersions}, \texttt{status} (draft, scheduled, published, completed), \texttt{shuffleConfig} (gồm shuffleQuestions và shuffleOptions), \texttt{paperEngine} (amc hoặc pdfkit).

\noindent\textbf{Bảng submissions.} Lưu trữ bài làm của học sinh với các trường chính: \texttt{\_id}, \texttt{examId}, \texttt{versionId} (tham chiếu phiên bản đề), \texttt{studentId}, \texttt{answers} (Array gồm questionId, selectedAnswer, isCorrect, score), \texttt{images} (danh sách URL ảnh phiếu), \texttt{scanMetadata} (thông tin quét), \texttt{omrSummary} (tóm tắt kết quả OMR), \texttt{status}, \texttt{manualOverrides}, \texttt{totalScore}, \texttt{timestamps}.

\noindent\textbf{Bảng appeals.} Lưu trữ đơn phúc khảo với các trường chính: \texttt{\_id}, \texttt{submissionId}, \texttt{examId}, \texttt{studentId}, \texttt{questionId}, \texttt{reason} (lý do), \texttt{expectedAnswer} (đáp án mong đợi), \texttt{evidenceImages} (danh sách URL ảnh minh chứng), \texttt{status} (pending, approved, rejected), \texttt{teacherResponse} (phản hồi giáo viên), \texttt{processedBy} (giáo viên xử lý), \texttt{processedAt}, \texttt{timestamps}.

Mối quan hệ giữa các bảng được thiết kế theo nguyên tắc tham chiếu thông qua ObjectId thay vì nhúng trực tiếp (trừ một số trường hợp đặc biệt như danh sách câu trả lời trong submission), giúp dễ dàng cập nhật và mở rộng. Các chỉ mục (index) được tạo trên các trường thường xuyên được truy vấn như email trong users, examId trong submissions, schoolId trong classes, nhằm tối ưu hiệu năng truy vấn.

\section{Xây dựng ứng dụng}
\label{section:4.3}

\subsection{Thư viện và công cụ sử dụng}
\label{subsection:4.3.1}

Danh sách các công cụ, ngôn ngữ, thư viện chính được sử dụng trong đồ án được tổng hợp trong Bảng \ref{table:4.1}.

\begin{table}[H]
\centering
\begin{tabular}{|p{2.8cm}|p{4cm}|p{2cm}|p{3.5cm}|p{3.5cm}|}
\hline
\textbf{Thành phần} & \textbf{Công cụ / Thư viện} & \textbf{Phiên bản} & \textbf{Mục đích} & \textbf{URL tải / Tài liệu} \\ \hline
Runtime & Node.js & 18.x trở lên & Môi trường chạy backend & \url{https://nodejs.org} \\ \hline
Framework backend & Express & 4.17.1 & Xây dựng REST API & \url{https://expressjs.com} \\ \hline
Cơ sở dữ liệu & MongoDB & 5.0 trở lên & Lưu trữ dữ liệu & \url{https://www.mongodb.com} \\ \hline
ODM & Mongoose & 5.7.7 & Ánh xạ đối tượng - tài liệu & \url{https://mongoosejs.com} \\ \hline
Xác thực & Passport.js + JWT & 0.4 / 9.0 & Xác thực và phân quyền & \url{https://www.passportjs.org} \\ \hline
Xác thực dữ liệu & Joi & 17.3.0 & Kiểm tra đầu vào & \url{https://joi.dev} \\ \hline
Upload file & Cloudinary SDK & 2.10 & Lưu trữ ảnh và PDF & \url{https://cloudinary.com} \\ \hline
AI & @google/generative-ai & 0.24.1 & Tích hợp Gemini & \url{https://ai.google.dev} \\ \hline
Email & Nodemailer & 6.3.1 & Gửi email thông báo & \url{https://nodemailer.com} \\ \hline
Kiểm thử backend & Jest & 26.x & Kiểm thử tự động & \url{https://jestjs.io} \\ \hline
Frontend & React & 18.x & Xây dựng giao diện web & \url{https://react.dev} \\ \hline
Ngôn ngữ web & TypeScript & 5.x & Kiểu tĩnh cho JavaScript & \url{https://www.typescriptlang.org} \\ \hline
Build web & Vite & 5.x & Công cụ đóng gói & \url{https://vitejs.dev} \\ \hline
State web & Zustand & 5.x & Quản lý trạng thái & \url{https://zustand-demo.pmnd.rs} \\ \hline
Di động & Flutter & 3.x & Đa nền tảng & \url{https://flutter.dev} \\ \hline
State di động & flutter\_bloc & 9.1.1 & Quản lý trạng thái BLoC & \url{https://bloclibrary.dev} \\ \hline
HTTP di động & Dio & 5.8.0 & Giao tiếp API & \url{https://pub.dev/packages/dio} \\ \hline
Camera di động & camera (Flutter) & 0.11.1 & Truy cập camera & \url{https://pub.dev/packages/camera} \\ \hline
OMR & opencv\_dart & 2.2.1 & Xử lý ảnh & \url{https://pub.dev/packages/opencv_dart} \\ \hline
IDE & Visual Studio Code & Latest & Soạn thảo mã nguồn & \url{https://code.visualstudio.com} \\ \hline
LaTeX & AMC (Auto-Multiple-Choice) & Latest & Sinh đề và phiếu OMR & \url{https://www.auto-multiple-choice.net} \\ \hline
XeLaTeX & TeX Live & 2023 & Biên dịch LaTeX & \url{https://tug.org/texlive} \\ \hline
WSL & WSL2 & Ubuntu 24.04 & Chạy LaTeX trên Windows & \url{https://learn.microsoft.com/windows/wsl} \\ \hline
\end{tabular}
\caption{Danh sách thư viện và công cụ sử dụng trong đồ án}
\label{table:4.1}
\end{table}

\subsection{Kết quả đạt được}
\label{subsection:4.3.2}

Sau quá trình phát triển, hệ thống Smart Grading đã được xây dựng hoàn chỉnh gồm các sản phẩm đóng gói:
\begin{itemize}
    \item backend API triển khai được trên Node.js, có thể chạy trực tiếp hoặc đóng gói Docker;
    \item giao diện web dạng SPA (Single Page Application) có thể build ra các file tĩnh và triển khai trên bất kỳ web server nào;
    \item ứng dụng di động Flutter có thể build ra file APK cho Android và IPA cho iOS.
\end{itemize}

Thống kê tổng quan về mã nguồn của hệ thống được tổng hợp trong Bảng \ref{table:4.2}.

\begin{table}[H]
\centering
\begin{tabular}{|p{3.5cm}|p{3cm}|p{3cm}|p{4cm}|}
\hline
\textbf{Thành phần} & \textbf{Số file mã nguồn} & \textbf{Số dòng mã (ước tính)} & \textbf{Ghi chú} \\ \hline
Backend (Node.js) & 190 & 17.000+ & Gồm services, controllers, models \\ \hline
Mobile (Flutter) & 199 & 28.000+ & Gồm screens, blocs, repositories \\ \hline
Web (React) & 151 & 21.000+ & Gồm pages, components, stores \\ \hline
Tổng cộng & 540+ & 66.000+ & Toàn bộ dự án \\ \hline
\end{tabular}
\caption{Thống kê mã nguồn của hệ thống Smart Grading}
\label{table:4.2}
\end{table}

Các lớp chính trong hệ thống gồm: 17 lớp model Mongoose tương ứng với 17 bảng dữ liệu; 16 service xử lý nghiệp vụ; 8 BLoC cho ứng dụng di động; 15 store Zustand cho giao diện web. Tổng dung lượng mã nguồn khoảng 4,5 MB, dung lượng ứng dụng di động sau khi build khoảng 25 MB cho Android.

\subsection{Minh họa các chức năng chính}
\label{subsection:4.3.3}

\noindent\textbf{Màn hình đăng nhập.} Người dùng truy cập hệ thống và đăng nhập bằng email và mật khẩu. Hệ thống hỗ trợ đăng nhập trên cả web và di động với giao diện thống nhất. Sau khi đăng nhập thành công, hệ thống chuyển hướng người dùng đến dashboard tương ứng với vai trò.

\noindent\textbf{Màn hình dashboard.} Dashboard hiển thị các chỉ số tổng quan phù hợp với vai trò người dùng: tổng số lớp, học sinh, đề thi, bài nộp, điểm trung bình, tỉ lệ đạt, danh sách bài nộp gần đây và số đơn phúc khảo đang chờ xử lý.

\noindent\textbf{Màn hình quản lý câu hỏi và đề thi.} Giáo viên có thể tạo câu hỏi thủ công hoặc sử dụng AI để sinh câu hỏi từ chủ đề. Màn hình tạo đề thi cho phép chọn câu hỏi từ ngân hàng, cấu hình tham số trộn và sinh nhiều phiên bản đề tự động.

\noindent\textbf{Màn hình quét phiếu trên di động.} Giáo viên sử dụng camera của thiết bị để chụp phiếu trả lời. Hệ thống tự động nhận dạng các ô tô và hiển thị kết quả chấm ngay trên ứng dụng với khả năng chỉnh sửa thủ công.

\noindent\textbf{Màn hình quản lý phúc khảo.} Giáo viên xem danh sách đơn phúc khảo, mở chi tiết để so sánh với đáp án và phản hồi cho học sinh. Học sinh theo dõi trạng thái đơn phúc khảo của mình trong mục "Phúc khảo của tôi".

\noindent\textbf{Màn hình báo cáo thống kê.} Hệ thống hiển thị các biểu đồ phổ điểm, tỉ lệ đạt theo lớp, độ khó và độ phân biệt của từng câu hỏi. Giáo viên có thể xuất báo cáo ra file PDF hoặc Excel để sử dụng ngoài hệ thống.

\section{Kiểm thử}
\label{section:4.4}

Hệ thống được kiểm thử theo hai kỹ thuật chính là kiểm thử hộp đen (black-box testing) và kiểm thử hộp trắng (white-box testing). Kiểm thử hộp trắng được áp dụng cho backend với framework Jest, bao gồm kiểm thử đơn vị (unit test) cho từng service và controller, kiểm thử tích hợp (integration test) cho luồng API end-to-end. Kiểm thử hộp đen được áp dụng cho cả web và di động thông qua kiểm thử thủ công và kiểm thử giao diện tự động với React Testing Library và Flutter Test.

Đồ án tập trung thiết kế các ca kiểm thử cho ba chức năng chính gồm:
\begin{itemize}
    \item Tạo đề thi và trộn câu cho nhiều đề thi;
    \item Chấm thi bằng ứng dụng;
    \item Phản hồi phúc khảo.
\end{itemize}
Tổng số ca kiểm thử được thiết kế cho ba chức năng này là 45 ca, trong đó có 18 ca cho chức năng tạo đề thi, 15 ca cho chức năng chấm thi và 12 ca cho chức năng phúc khảo. Kết quả kiểm thử cho thấy 42/45 ca đạt kết quả như mong đợi, 3 ca không đạt đã được phân tích và sửa lỗi. Bảng \ref{table:4.3} tổng hợp kết quả kiểm thử tổng quan.

\begin{table}[H]
\centering
\begin{tabular}{|p{4cm}|p{3cm}|p{3cm}|p{3cm}|}
\hline
\textbf{Chức năng} & \textbf{Số ca kiểm thử} & \textbf{Số ca đạt} & \textbf{Tỉ lệ đạt} \\ \hline
Tạo đề thi, trộn câu & 18 & 17 & 94.4\% \\ \hline
Chấm thi bằng ứng dụng & 15 & 13 & 86.7\% \\ \hline
Phản hồi phúc khảo & 12 & 12 & 100\% \\ \hline
\textbf{Tổng cộng} & \textbf{45} & \textbf{42} & \textbf{93.3\%} \\ \hline
\end{tabular}
\caption{Tổng hợp kết quả kiểm thử các chức năng chính}
\label{table:4.3}
\end{table}

Ba ca kiểm thử không đạt tập trung ở chức năng chấm thi bằng ứng dụng, nguyên nhân chính là do ảnh chụp trong điều kiện ánh sáng yếu hoặc bị nghiêng quá nhiều. Sau khi điều chỉnh thuật toán tiền xử lý ảnh với bước tăng cường độ tương phản thích ứng và bước chỉnh nghiêng cải tiến, các ca này đã đạt trong lần kiểm thử lại.

\section{Triển khai}
\label{section:4.5}

Hệ thống Smart Grading đã được triển khai thử nghiệm trong môi trường thực tế tại hai lớp học với tổng cộng khoảng 60 học sinh trong vòng 4 tuần. Môi trường triển khai gồm: máy chủ backend chạy trên Node.js với cấu hình 4 CPU, 8 GB RAM, ổ cứng SSD 100 GB; cơ sở dữ liệu MongoDB chạy trên cùng máy chủ với ổ cứng riêng; dịch vụ Cloudinary dùng để lưu trữ ảnh và PDF; giao diện web được triển khai trên nginx; ứng dụng di động được cài đặt trên điện thoại Android của giáo viên.

Trong quá trình thử nghiệm, hệ thống đã được sử dụng để tổ chức 5 bài kiểm tra trắc nghiệm với tổng số khoảng 300 bài làm được chấm. Kết quả thống kê cho thấy:
\begin{itemize}
    \item thời gian chấm trung bình cho một phiếu trên ứng dụng di động là 2,3 giây;
    \item độ chính xác trung bình của thuật toán nhận dạng OMR đạt 96,2\%;
    \item thời gian phản hồi API trung bình là 180ms;
    \item hệ thống xử lý được đồng thời 50 người dùng truy cập mà không suy giảm hiệu năng đáng kể.
\end{itemize}
Phản hồi từ giáo viên và học sinh tham gia thử nghiệm đều tích cực, đặc biệt về tính tiện lợi của việc chấm bài trên di động và sự minh bạch trong quy trình phúc khảo.

\section{Kết chương}
Chương 4 đã trình bày toàn bộ quá trình phân tích thiết kế, triển khai và đánh giá hệ thống Smart Grading. Trong phần thiết kế kiến trúc, chương đã lựa chọn mô hình phân lớp kết hợp client--server, mô tả biểu đồ gói UML với sự phụ thuộc một chiều giữa các tầng và biểu đồ lớp chi tiết cho nhóm gói nghiệp vụ chính. Trong phần thiết kế chi tiết, chương đã đặc tả hai lớp chủ đạo \texttt{ExamService} và \texttt{SubmissionService}, mô tả luồng truyền thông điệp cho hai use case quan trọng, đồng thời trình bày thiết kế giao diện với chuẩn phối màu và kích thước cụ thể cùng mô tả bốn bảng cơ sở dữ liệu cốt lõi. Phần xây dựng ứng dụng đã thống kê danh sách thư viện với phiên bản chính xác, các thông số mã nguồn (hơn 540 file, 66.000 dòng mã), kèm minh họa các chức năng chính. Phần kiểm thử đã thiết kế 45 ca cho ba chức năng chính với tỉ lệ đạt 93,3\%. Phần triển khai thử nghiệm tại hai lớp học với 60 học sinh trong bốn tuần cho thấy hệ thống đáp ứng được các yêu cầu phi chức năng đã đặt ra. Chương 5 sẽ tập trung trình bày chi tiết các giải pháp kỹ thuật nổi bật và đóng góp chính của đồ án.

\end{document}